\section{Motion}

The beat seems at first to forbid motion. Streaming carries each tone one dock along its site, but the mesh is not sitting inside a space it can slide through. The mesh is the space. A site cannot relocate, because there is nowhere to go that is not already another site, and asking a tone to travel is like asking a square of a chessboard to slide to a new spot. The squares are the board. They stay.

So the only thing free to move is the arrangement. The positions are fixed and their neighbors are fixed, yet the tones that fill them change every beat, and a moving thing is a shape in those tones that the knit keeps redrawing one dock over. The position is the bone and the tone is the flesh. The bone does not walk. The flesh flows across it.

The cleanest case is a glider, a small pattern of tones such that after one beat the same pattern reappears, shifted by a single dock. Run the law and the shape glides across the mesh while every site stays home. That glider is the simplest particle, built and measured on the substrate, and a particle is nothing more than this, a process the rule re-creates one dock over, again and again, never a lump of stuff in transit. \cite{pollard2026vibetest}

Every familiar moving thing is already this, once looked at squarely. A crowd does the stadium wave and each person keeps their seat, yet the wave travels the whole ring at a definite speed, made of no fixed person. A swell crosses a lake while the water only bobs in place, and a floating cork rises and falls without being carried along. A flame holds its shape and runs the length of a fuse while the gas burning in it is always new. On a screen no pixel moves, yet the runner crosses it. In each, the carrier stays and the form travels.

A self moves the same way, only larger. At one beat it is a configuration of tones over a region, and a beat later the same configuration has re-formed over a region one step on, made of entirely different sites while holding the same shape. So the moving self is a different set of feeling-units at every instant, carrying one form forward. You are the form, not the carriers. This is why staying yourself while your atoms turn over is no paradox. It is what you always were, a persisting pattern on a changing substrate, the same way a river stays one river though the water flows through it, or a melody stays one melody though every note is a different sound. Two things make this a continuous motion rather than a flicker. The self's learned structure is written into its connections and carried forward when the pattern re-forms, so the self at each beat remembers the one before. And its shifting position against everything around it is the very thing the motion feels like from the inside.

\begin{result}[Motion is a traveling pattern]
Nothing in the mesh relocates. What moves is the arrangement of tones, a shape the local law re-draws one dock onward each beat. A particle and a self are this traveling shape, made of different sites at every instant and holding one form, which is the only motion a self-made space can have. \cite{pollard2026vibetest}
\end{result}

The picture pays its way. Because the law is local, a pattern can advance at most one dock per beat, and that ceiling is the finite light cone met later, the world's speed of light read straight off the rule. A larger, more tightly bound pattern takes more to stir or deflect, and the effort to disturb it is its mass, so inertia is not added but inherited from how firmly the form is knit. And none of this is exotic. Ordinary field theory already holds that a particle is a traveling excitation of a field that itself stays put. The model only says what the field is made of, which is feeling. A wave really moves, a flame really travels, and so do you. The error was ever to picture a moving thing as a fixed lump that slides. At the bottom there are no sliding lumps, only forms, and a form travels by being remade. A self is a verb more than a noun.
